\justifying
\NSFCBodyText

% \subsubsection{年度研究计划}

% 表\ref{tab:annual_plan}展示了本项目三年期的详细研究计划与预期成果。

% \begin{table}[htbp]
%   \centering
%   \fontsize{10}{10}\selectfont
%   \caption{\textbf{三年期研究计划与预期成果}}
%   \begin{tabular}{lccc}
%   \toprule
%   \textbf{研究阶段} & \textbf{研究重点} & \textbf{主要任务} & \textbf{预期成果} \\
%   \midrule
%   第一年 & 数据收集与建模 & 建立数据库、文献综述 & 论文1-2篇、理论框架 \\
%   第二年 & 深度分析与验证 & 序列分析、形象分析 & 论文2-3篇、模型原型 \\
%   第三年 & 模型验证与转化 & 案例比较、工具开发 & 论文3-4篇、专著、工具 \\
%   \bottomrule
%   \end{tabular}
%   \captionsetup{font={footnotesize,stretch=1.25},justification=raggedright}
%   \label{tab:annual_plan}
% \end{table}

% \subsubsubsection{第一年研究计划}

% \noindent\textbf{研究重点}

% 数据收集与预处理、职业轨迹初步建模。

% \noindent\textbf{主要工作}

% 完成佐佐木希职业发展档案数据库建立（媒体报道、影视作品、商业代言、杂志封面、社交媒体舆情数据），开发数据清洗与预处理流程。完成相关研究文献综述，构建“艺人职业发展三阶段动态模型”理论框架，开发“艺人职业轨迹编码手册”初稿。

% \noindent\textbf{年度预期成果}

% 建立完整研究数据库；完成文献综述论文1-2篇；构建理论模型框架。

% \subsubsubsection{第二年研究计划}

% \noindent\textbf{研究重点}

% 深度数据分析、模型构建与验证。

% \noindent\textbf{主要工作}

% 应用序列分析方法识别职业发展关键转型节点与典型路径，构建“职业状态转移矩阵”，开发“职业转型预测模型”原型。采用深度学习方法对影像资料进行特征提取，运用NLP技术分析文本资料中的形象建构策略，构建“跨媒介形象一致性测量指标体系”。

% \noindent\textbf{年度预期成果}

% 完成职业轨迹建模分析；发表实证研究论文2-3篇；开发模型原型系统。

% \subsubsubsection{第三年研究计划}

% \noindent\textbf{研究重点}

% 模型验证、案例比较、成果总结与转化。

% \noindent\textbf{主要工作}

% 通过交叉验证和敏感性分析评估模型稳健性，选取同时期其他日本女性艺人作为对照案例检验外部效度。撰写研究总报告，开发“艺人职业规划决策支持工具”原型，组织学术研讨会，撰写专著《艺人职业发展的动态机制：基于佐佐木希案例的纵向研究》。

% \noindent\textbf{年度预期成果}

% 完成模型验证与优化；发表高水平研究论文3-4篇；完成研究总报告；开发决策支持工具原型。

% \subsubsubsection{预期研究结果}

% \noindent\textbf{理论贡献}

% 提出“艺人职业发展三阶段动态模型”；建立“品牌韧性”理论框架；构建“跨媒介形象一致性测量指标体系”；开发“艺人职业轨迹编码手册”；提出多模态数据融合分析框架。

% \noindent\textbf{实践价值}

% 为经纪公司提供“艺人职业规划决策支持工具”；为艺人本人提供“个人品牌管理指南”；为政府文化部门制定政策提供实证依据。

% \noindent\textbf{学术产出}

% 计划发表高水平学术论文6-8篇（SSCI/CSSCI不少于4篇），目标期刊包括 Journal of Business Research、Tourism Management、《新闻与传播研究》《现代传播》等；撰写专著1部；参加国际学术会议2-3次；组织专题学术研讨会1-2次。


%--------------------------------------------------------------------------
% 第三部分：年度研究计划（3年期）
%--------------------------------------------------------------------------
\subsubsection{年度研究计划及预期研究结果}

\subsubsubsection{年度研究计划（2027.01 --- 2029.12）}

本项目研究周期为三年，严格遵循“多维感知解耦 $\rightarrow$ 核心架构研发 $\rightarrow$ 物理迁移部署”的层级递进科学逻辑。表~\ref{tab:annual_plan} 展示了本项目每年的阶段划分、研究重点、核心任务与预期产出。具体而言，各年度的实施方案与执行细节安排如下：

% 插入年度研究计划汇总表
\begin{table}[htbp]
  \centering
  \fontsize{10}{13}\selectfont % 调整字号与行距，提升学术排版质感
  \caption{\textbf{本项目三年期研究计划与预期成果汇总表}}
  \label{tab:annual_plan}
  \begin{tabular}{p{0.1\textwidth} p{0.22\textwidth} p{0.38\textwidth} p{0.2\textwidth}}
  \toprule
  \textbf{研究阶段} & \textbf{研究重点} & \textbf{核心任务} & \textbf{预期成果} \\
  \midrule
  \textbf{第一年度} \newline (2027年) & 多模态感知解耦与\newline 物理闭环映射 & 基于Analysis-by-Synthesis的时空同步解析；构建融合Isaac Lab动力学反馈的隐空间重定向机制 & 构筑并开源 Embodied-Social-100k 数据集；\newline 论文1篇，专利1项 \\
  \midrule
  \textbf{第二年度} \newline (2028年) & Social-MLA核心架构\newline 研发与生成演化 & 研发多模交互Token预测（MITP）模块；揭示离散语义对常微分方程（ODE）连续演化轨迹的调制机理 & 确立Social-MLA算法原型；\newline 论文1篇、专利1项 \\
  \midrule
  \textbf{第三年度} \newline (2029年) & 虚实迁移控制策略\newline 与实景验证部署 & 设计融合底盘平衡的RL抗扰动追踪策略；在履带式人形机器人上开展高表现力交互的主客观综合测评 & 完整软硬件验证系统；\newline 论文1篇、专利1项 \\
  \bottomrule
  \end{tabular}
  \captionsetup{font={footnotesize,stretch=1.25},justification=raggedright}
\end{table}


\paragraph{\textbf{第一年度（2027.01 --- 2027.12）：复杂交互感知解耦与物理闭环数据集构筑}}
本年度聚焦于打通数据壁垒，为生成大模型提供高质量的感知先验与数据基座。
\begin{itemize}
    \item \textbf{机理探究与算法研发：} 重点研发基于音画一致性的高价值交互数据自动化筛选管线；构建基于Analysis-by-Synthesis范式的闭环时空解析网络，突破遮挡环境下“表情-视线-动作”的联合解耦与高保真参数化重构难题。
    \item \textbf{仿真搭建与动力学校验：} 搭建目标履带式人形机器人的Isaac Lab高保真物理仿真环境并精确标定动力学参数；探究隐空间动作重定向网络，将物理引擎的力矩与质心偏移惩罚作为梯度反馈注入网络，完成动力学边界补偿机制的训练。
    \item \textbf{阶段性测试：} 跑通“Video-to-Sim”自动化数据流转闭环，完成至少10万规模级的高表现力对齐样本生成，正式定型 Embodied-Social-100k 数据集。
\end{itemize}

\paragraph{\textbf{第二年度（2028.01 --- 2028.12）：Social-MLA 核心生成架构设计与跨模态对齐研究}}
本年度为本项目的核心攻坚期，全面开展基于大模型与流匹配的跨模态协同生成机制研究。
\begin{itemize}
    \item \textbf{架构搭建与级联自回归设计：} 设计异构多模态特征编码器与融合模块；研发多模交互Token预测（MITP）机制，利用大语言模型完成离散语义的自回归时序规划，并在解码路径中引入渐进更新的上下文条件信息流。
    \item \textbf{连续流形演化与模型联调：} 深度探究流匹配模型在具身动作生成中的数学表达，推导并实现离散语义Token对连续动作常微分方程（ODE）非线性漂移项的条件调制；利用第一年构筑的数据集对 Social-MLA 架构进行端到端联合训练。
    \item \textbf{阶段性测试：} 建立离线评测基准，利用语义一致性分数、动作平滑度等客观指标，论证架构在跨模态对齐与时空协同生成上的优越性。
\end{itemize}

\paragraph{\textbf{第三年度（2029.01 --- 2029.12）：底层鲁棒控制策略研制与适老化真机场景落地}}
本年度重点攻克虚实鸿沟（Sim-to-Real），验证项目理论在真实物理世界中的系统效能。
\begin{itemize}
    \item \textbf{追踪策略与抗扰动研究：} 针对履带式机器人构型，研究基于强化学习（RL）的底层追踪控制策略；设计融合底盘稳定性约束（如俯仰/横滚角惩罚）与动作跟踪精度的复合奖励函数，探索域随机化技术以克服Sim-to-Real摩擦与延迟分布差异。
    \item \textbf{系统集成与实景验证：} 将训练收敛的 Social-MLA 生成模型与 RL 追踪策略无缝部署至真实履带式双臂人形机器人硬件平台；在适老化陪护、自然语音问答等典型交互场景中，开展大规模的人机交互主观实验与底盘抗扰动压力测试。
    \item \textbf{项目结题与成果总结：} 整理实验数据，撰写高水平学术论文与项目结题报告，推动相关数据集与核心算法的开源共享工作。
\end{itemize}

\subsubsubsection{预期研究结果}

本项目预期在理论机理突破、核心算法研发、数据生态构建以及人才梯队培养等方面取得如下实质性成果：

\begin{enumerate}
    \item \textbf{基础理论与核心算法：} 提出并完善一套面向共身智能的跨模态动作感知、映射与生成理论框架；从底层数学机理上揭示离散意图Token驱动连续具身动作ODE流形演化的跨模态调制规律；构建并定型 Social-MLA 高表现力动作生成算法原型。
    \item \textbf{高保真数据集与开源生态：} 构筑并开源国际上首个兼顾动力学约束与多模态情感语义严格对齐的高保真具身社交数据集 \textbf{Embodied-Social-100k}；开源相关的“分析-合成”多维解析代码与物理闭环映射接口，填补学术界在高质量共身交互数据上的空白。
    \item \textbf{知识产权与高水平学术论文：} 预期在 \textit{IEEE TPAMI, IJCV, CVPR, ICCV, ICLR, ICRA} 等人工智能、计算机视觉与机器人领域内顶级国际期刊与会议上发表高质量学术论文 3-4 篇；针对Social-MLA核心生成架构与底盘平衡追踪策略，申请并获得授权国家发明专利 3-4 项。
    \item \textbf{物理验证系统与人才梯队建设：} 形成一套软硬件协同的高表现力履带式机器人社会化交互验证系统；以本项目的实施为重要契机，培养 2-3 名精通图形学生成算法与机器人底层控制的交叉学科拔尖研究生，构筑具有强竞争力的青年科研创新团队。
\end{enumerate}